\section{Discussion}
\label{sec:discussion}

The primary finding is the gap between generation availability and executable test quality. GPT returned content for every target, yet 49 suites remained invalid after the allowed repair step. This makes the full-set branch coverage and mutation score modest, and it prevents the predefined RQ1, RQ2, and RQ4 hypotheses from being supported. The result is useful even though it is negative: in this protocol, zero-shot generation did not consistently produce valid test oracles for the chosen Java targets.

The paired results require a different interpretation. When both GPT and EvoSuite suites pass, the selected GPT suites have high coverage and mutation scores, and the five-minute mutation comparison has the smallest raw p-value. However, the paired sample contains only 14 executable GPT suites (13 for mutation) and the branch analysis has two non-zero ranked differences. After controlling the six RQ3 comparisons with Holm adjustment, no difference is significant. These results do not establish parity between tools; they establish that this run lacks statistically significant evidence of a difference on the small pass-conditioned subset.

For practice, the observed workflow suggests that an LLM can be a useful candidate-test generator only when followed by strict compilation, execution, and oracle validation. A raw response should not be counted as a successful test suite. Search-based generation remained consistently executable on the archived EvoSuite suites under the measured budgets, while the LLM pipeline's main improvement opportunity is repair or feedback targeted at assertion validity.

The study also clarifies the scope of the comparator. A separate feasibility evaluation carried out prior to this run confirmed that EvoSuite branch coverage reaches a plateau at 99.61\% for budgets beyond 30 minutes on this corpus \cite{Lemieux_2023}; the three remaining uncovered branches are structurally unreachable by the search regardless of budget. The 1-, 3-, and 5-minute budgets used in RQ3 therefore span a meaningful range of EvoSuite performance---rising from 90.29\% to 99.34\% branch coverage---while stopping well below the point of diminishing returns. EvoSuite supplies a reproducible engineering baseline on the same SUTs; it does not measure what students write. A future study should collect student-written suites and evaluate them with the identical build, coverage, and mutation pipeline before making any GPT--student claim.

\subsection{Practical Implications}
For a developer or teaching team, the immediate implication is not to reject LLM assistance altogether. Instead, the generation step should be treated as a proposal mechanism. A safe workflow keeps the generated source, compiles it in isolation, executes it against the known-correct SUT, and measures it only after it passes. A failure category should remain visible in the final dashboard. This workflow prevents a high-level language-model success rate from obscuring the cost of manual debugging or incorrectly asserted behavior.

The result also motivates feedback-guided repair. The present run allowed only one repair invocation and did not use coverage or mutation feedback to steer subsequent generation. Because assertion failures dominate the observed invalid outcomes, a next protocol can give the model the compiler output, failing assertion, or minimized counterexample and then evaluate whether a bounded repair policy raises the executable-suite rate without inflating API cost \cite{yang2025react, pham2025llmloop}. Such a follow-up should predefine the maximum repair attempts and report the separate contribution of each attempt.

\subsection{Interpretation of Negative Findings}
The negative threshold results should not be described as evidence that LLMs cannot generate useful tests. They establish a narrower fact: this model snapshot, fixed prompt, and zero-shot repair policy did not meet the preregistered thresholds on this 63-SUT setting. The evidence is informative because it identifies which condition failed most often. A future comparison can alter one factor at a time--for example, structured few-shot prompting, test-feedback loops, or domain-specific context \cite{vy2025whatinputs}--while retaining the same execution and metric gates.

\subsection{Measurement and Reporting Implications}
This evaluation also demonstrates why a single aggregate number is insufficient for an LLM testing study. The all-target scores answer a deployment-oriented question: if a team runs the configured process once for every selected function, what quality is obtained without discarding failures? The pass-conditioned paired analysis answers a different diagnostic question: when both approaches produce runnable suites for the same target, is there evidence that their measured adequacy differs? Mixing the denominators would hide the practical failure rate in the former question or unfairly attribute failures to a tool in the latter question. The report therefore keeps the 63-target population, the executable GPT subset, the paired subset, and the non-zero ranked subset explicit in both prose and tables.

The use of mutation analysis has a similar implication. Branch coverage indicates whether tests traverse decisions, whereas mutation score asks whether their assertions expose seeded behavioral changes. A generated suite can execute many branches while still accepting incorrect values, especially when assertions are weak, over-general, or based on an incorrect interpretation of the specification. Conversely, a small test suite may kill a high proportion of the mutants reachable in a simple function. Reporting the two measures together does not make them interchangeable; it makes their distinct limits visible. For this reason, future extensions should retain the class-level counts behind each percentage and should classify the major invalid-suite modes, rather than only reporting a mean coverage score.

Finally, the study treats reproducibility as part of the result rather than a later packaging task. The retained generation records make it possible to rerun the computations, inspect a particular invalid class, and distinguish a model-output problem from a build or instrumentation problem. This is particularly important for API-based experiments because model aliases, service behavior, and local dependencies can change over time. A replication should preserve the dated model identifier, prompt template, retry and repair policy, target manifest, compiler version, metric-tool versions, and exact decision rules before comparing a new run with this baseline.


\subsection{Operational Cost and Selection Effects}
The recorded service cost is small for this run, but cost should not be interpreted independently of validity. The initial calls cost \$0.030982 and the permitted repair calls cost \$0.036773. These figures describe only the API portion of the workflow. They omit the local time needed to compile, diagnose a failure, measure the suite, and decide whether further human intervention is acceptable. In a practical setting, a comparison between test-generation methods should therefore report at least three dimensions: generation expenditure, executable-suite yield, and post-generation engineering effort. A cheap response that fails its assertions may still be expensive when it creates review work for a developer.

The RQ3 subset creates a second, statistical form of selection. By definition, it excludes the 49 GPT suites that did not execute. This is appropriate for a paired measurement of score differences because coverage and mutation scores cannot be compared fairly when one side has no valid test class. It is not appropriate as a stand-alone claim about end-to-end test generation. The full-set and paired views should consequently be read together. The former bounds the outcome a user receives from the configured pipeline; the latter describes the quality of the successful cases. Presenting both avoids the common but misleading practice of reporting high coverage only after silently filtering failed generations.

\subsection{Design for a Follow-Up Study}
A direct follow-up should change as few variables as possible while targeting the observed bottleneck. One straightforward design holds the 63-SUT corpus and the JaCoCo/PIT measurement configuration constant and varies the prompt strategy: fixed zero-shot, structured few-shot, and bounded feedback-guided repair. Each condition should use the same dated model identifier, the same maximum repair attempts, the same temperature, and the same token-accounting policy. The primary reported outcome should remain the executable-suite rate, with branch and mutation scores reported separately for the full corpus and the passing subset. Such a design would directly answer whether additional context or iterative feedback improves the rate of valid oracle construction under otherwise identical experimental conditions.
